\documentclass[../main.tex]{subfiles}
\begin{document}
%==========================================TIÊU ĐỀ TẬP========================================
\begin{table}[h]
    \begin{adjustwidth}{-1cm}{}
        \begin{tabular}{>{\centering\arraybackslash}p{6cm}|>{\raggedright\arraybackslash}p{12cm}}
            \multirow{5}{6cm}
            % Cột bên trái: ÉPISODE và số tập
            {\\[-40pt]
                \centering
                {\fontsize{20pt}{20pt}\selectfont{\outlinetext[1]{eptype}{white}{\flaregothic ÉPISODE}}\\[12pt]
                    {\fontsize{72pt}{72pt}\selectfont{\outlinetext[1]{epnum}{white}{\burbank 163}}}\\[6pt]       % use for normal/special episodes
                    % {\fontsize{36pt}{36pt}\selectfont{\outlinetext[1]{epnum}{white}{\burbank SÁU}}}\\[6pt]     % use for custom episodes
                    {\fontsize{20pt}{20pt}\selectfont{\outlinetext[1]{epnum}{white}{\cafeteria (13-4)}}}\color{black}}
                \\[18pt]
                {\fontsize{14pt}{14pt}\selectfont{\outlinetext[1]{epattr}{white}{\flaregothic BIENVENUE, }}{\flaregothic\color{Gozaru} Iroha}\fontsize{14pt}{14pt}\selectfont{\outlinetext[1]{epattr}{white}{\flaregothic \ !}}}\\
                % {\fontsize{18pt}{18pt}\selectfont{\outlinetext[1]{epattr}{white}{\flaregothic [II]}}}\\
                \color{black}
            }
             &
            % Dòng thứ nhất: tên tập tiếng Nhật
            {\fotskip\fontsize{18pt}{18pt}\selectfont \ruby{先}{せん}\ruby{輩}{ぱい}に\ruby{喧}{けん}\ruby{嘩}{か}\ruby{売}{う}ったらボコボコにされた\ruby{話}{はなし}}              \\[6pt]
            % Dòng thứ hai: tên tập tiếng Anh
             & {\cafeteria\fontsize{18pt}{18pt}\selectfont Choose Your Battles Wisely} \\[4pt]
            % Dòng thứ ba: tên tập tiếng Việt
             & {\vnmsans\fontsize{18pt}{18pt}\selectfont Thách đấu với kẻ mạnh nhất}   \\[8pt]
            % Dòng thứ tư: Nhân vật xuất hiện
             & {\caviar{\fontsize{14pt}{14pt}\selectfont Nhân vật: \emp{kuon2} \emp{ryuuhou2} \emp{game} \emp{sora1} \emp{roboco2} \emp{haato} \emp{ayame} \emp{noel} \emp{iroha}}} \\[8pt]       
            %  \\
            % Dòng thứ năm (nếu có): Nhân vật chỉ xuất hiện trong Omake
            %  & {\abv Truyện phụ:} \emp{okayu}
            \\[8pt]
            % Dòng cuối cùng: Thời gian diễn ra của tập (thường sẽ là ngày phát hành)
             & {\sen\fontsize{16pt}{16pt}\selectfont Mốc thời gian: 26/06/2022}        \\[18pt]
        \end{tabular}\\[8pt]
    \end{adjustwidth}
\end{table}
%=========================================TRUYỆN CHÍNH========================================
\color{black}

\txtbx[2]{
    {\color{vol} \uline{Tóm tắt:}} Kazama Iroha - thành viên mới của holoX - vừa đến văn phòng đã thách đấu Ayame. Trận đấu diễn ra chóng vánh: Iroha thắng dễ dàng, thanh kiếm của Ayame gãy làm đôi. Ayame hoảng loạn che giấu thất bại. Về văn phòng, Haato lao vào tấn công Iroha, rồi Roboco (bắn laser) và Noel (thây ma đói) truy đuổi. Sora xuất hiện, dẹp loạn bằng uy quyền. Iroha ngất xỉu vì sợ ``sát khí'' của Sora. Mọi thứ hóa ra là do ``Sếp'' của cô đã ra lệnh cho Iroha đi khiêu chiến tất cả mọi người, trừ Kuon. Kuon được đồn là ``kẻ mạnh nhất'' công ty, khiến Iroha từ nay quỳ lạy van xin tha mạng.
}
\pov{Hikawa Kuon}

*thở dài* Chậc. Lại chuẩn bị có thêm một Thế hệ nữa xuất hiện.

Tôi vừa mới sống sót qua chuỗi ngày sóng gió và hàng tá trò nghịch dại không có điểm dừng của các cô gái NePoLaBo - từ những màn xiếc bay lượn ngã cống của Polka, trò bắn tỉa súng hơi cay của Botan, nết nhậu nhẹt bí tỉ của Lamy, cho đến cả danh sách hàng triệu ``người chồng'' không hồi kết của Nene.

Tôi đã ngây thơ nghĩ rằng sau khi thế hệ thứ Năm ổn định, mọi thứ có lẽ sẽ dần đi vào quỹ đạo hơn, dù gì tôi cũng đã quen với việc các cô gái trong công ty cứ nhốn nháo như chốn không người rồi. Nhưng không, COVER dường như chưa bao giờ có khái niệm ``bình yên'' trong từ điển nhân sự. Tôi thực sự nên nộp đơn đòi chú Tanigo tăng gấp đôi lương, hoặc chí ít là trợ cấp thêm tiền mua thuốc hạ sốt và cà phê chống đau đầu.

Bởi vì ngay tuần này, tôi lại phải chuẩn bị tinh thần để đón tiếp năm thành viên của một ``tổ chức bí mật'' vừa mới đặt chân vào công ty - \textit{Secret Society holoX}, hay còn gọi là Thế hệ thứ Sáu. Nghe đâu bọn họ tự xưng là một tổ chức hắc ám ngầm với tham vọng xâm lăng và thống trị toàn bộ thế giới loài người\dots\ Nghe thì đao to búa lớn và nguy hiểm như trùm cuối trong phim siêu anh hùng vậy đấy, nhưng tôi dám cá một chầu ramen là cái gọi là ``tổ chức tà ác'' này sớm muộn gì cũng bị cái văn phòng tấu hài này đồng hóa thành một gánh xiếc mới mà thôi.

Ít nhất thì, tôi không còn phải chinh chiến với các ``lính mới'' này đơn độc. Vì cùng với tôi là các em gái sẽ đồng hành với tôi trong suốt năm tuần tiếp theo này.

Vậy chúng ta có gì cho tuần đầu tiên nào?

Một người con gái gắn bó với kiếm thuật, đến từ một vùng đất xa xôi, luôn nuôi dưỡng tinh thần võ sĩ đạo trong lòng và ấp ủ giấc mơ trở thành một ninja thực thụ.

Em ấy có một nghị lực thật đáng ngưỡng mộ; mặc dù ngoại hình khá mảnh mai, thế nhưng đừng để vẻ ngoài đó khiến bạn xem nhẹ năng lực mà em sở hữu. Với trình độ kiếm thuật được mài giũa qua vô số luyện tập, cùng với sự kỷ luật bản thân và sự kiên định không gì lay chuyển được, em luôn sẵn sàng đương đầu với mọi khó khăn, chỉ để che chở cho những người mà em đặt trọn niềm tin và yêu thương.

Thế nhưng, ẩn sau vẻ mặt lúc nào cũng chững chạc, nghiêm túc kia lại là một cá tính vô cùng đáng yêu: em rất lịch sự, biết lắng nghe và thường có chút ngại ngùng khi tiếp xúc với người mới; đặc biệt luôn thể hiện lòng kính trọng sâu sắc với những người em xem là người lãnh đạo của mình. Dù có đôi lần em hơi vụng về trong cách giao tiếp và dễ bị rơi vào trạng thái bối rối, em vẫn luôn nỗ lực hết mình để hoàn thành tốt mọi nhiệm vụ được giao. Một kiếm sĩ luôn giữ thái độ nghiêm túc trong công việc, thế nhưng cũng có những lúc em tỏ ra ngây thơ, non nớt chẳng khác nào một chú ngựa con mới chập chững bước vào đời.

Không những vậy, em còn là một bậc thầy trong việc\dots\ chinh phục những chặng đường dài bằng xe đạp và dành tình yêu thương để chăm sóc các loài động vật.

Và mặc dù em đã là một kiếm sĩ tài năng, ước mơ của em không chỉ dừng lại ở việc trở thành một ninja xuất sắc. Em mong rằng một ngày không xa, em sẽ trở thành một người mà ai cũng có thể tin tưởng giao phó, đồng thời cùng cả tổ chức holoX đạt được những thành tựu càng lúc càng lớn mạnh hơn.

Một nữ kiếm sĩ với lòng trung thành son sắt\dots

\dots hoặc có lẽ, một ``ninja tập sự'' vẫn đang trên hành trình tìm kiếm con đường của riêng mình. À, không, là một samurai mới đúng.

\chrint

Kazama Iroha (\ruby{風}{かざ}\ruby{真}{ま}いろは, \ruby{Ca-da}{Phong}-\ruby{ma}{Chân}\ I-rô-ha) - thành viên đầu tiên của tổ chức bí mật Secret Society holoX, hay còn gọi là Thế hệ thứ Sáu của \textit{hololive}. Nhiều người thường gọi cô là ``chiếc bảo hiểm đi động'' của holoX, và quả thật không có từ nào để diễn tả vai trò của cô nàng chính xác hơn thế.

Sinh ra và lớn lên ở một vùng núi hẻo lánh xa xôi, cô ấy đã rời bỏ quê hương cùng người bạn đồng hành của mình để bắt đầu một hành trình dài, với ước mơ được khám phá và học hỏi thêm bao điều kỳ diệu từ thế giới bên ngoài. Ngày nay, Iroha đảm nhận mọi nhiệm vụ bảo vệ an ninh và làm vệ sĩ cho các thành viên khác trong holoX, vừa đủ để trang trải cuộc sống hàng ngày.

Iroha là một cô gái nói chuyện nhẹ nhàng, tinh nghịch, nhưng đôi lúc cũng hơi vụng về một chút. Nhiều người hâm mộ nhận xét cô là thành viên ``trong sáng, đoan trang'' nhất của thế hệ mình, và cũng là một trong những người giữ được phong thái đó nhất trong toàn bộ \hll, ngang hàng với những cái tên như Tokino Sora - dù bản thân thì không bao giờ tự nhận mình là như vậy. Có lúc người ta còn đùa gọi cpp là một cô gái ``cơ bắp mà não'', vì trong buổi họp mặt đầu tiên của holoX, cô là người duy nhất không dám tự nhận mình là thành viên thông minh nhất trong nhóm.

\model{26}{r}{9cm}{iroha}

Còn có một điều đặc trưng không thể nhầm lẫn của Iroha chính là lối nói chuyện cổ xưa, với cụm từ ``de gozaru'' thường trực ở cuối gần như mỗi câu nói của cô, tạo nên nét cổ điển, trang trọng đúng như một võ sĩ thời phong kiến. Bản thân Iroha luôn tự hào nhận mình là một samurai chính hiệu, luôn trung thành và sẵn sàng bảo vệ đồng đội, nhưng thường xuyên lại bị các thành viên khác trong nhóm và fan hâm mộ trêu chọc, gắn cho cái mác ``ninja'' - tới mức độ mà cô nàng phải phồng má lên. Cô nàng thực sự không thích bị gọi như vậy.

Về ngoại hình, Iroha sở hữu mái tóc vàng ngắn được buộc thành một cái đuôi ngựa bằng dải băng hình lá cây, cùng với đôi mắt xanh lục trong veo như bầu trời ngày nắng. Trang phục thường ngày của cô nàng khá đặc biệt: một chiếc áo bolero ôm sát người, ngắn tay, với những băng vải quấn ngực thay vì áo trên bình thường; cô cũng có một chiếc áo khoác trắng với chi tiết xanh-vàng, được cài vào chiếc thắt lưng màu be và xanh lam quanh eo, từ đó buông xuống một sợi dây đỏ. Cô mặc quần short trong mỏng manh, bên ngoài là chiếc chân váy màu xanh ngọc với những họa tiết đậm hơn, đi kèm đôi găng tay đen lỡ ngón. Phía trên cùng là một chiếc áo khoác dáng giống kimono, cùng tông màu vàng với mái tóc, được đính thêm các chi tiết xanh và lá cây. 

Và dĩ nhiên, không thể thiếu thanh kiếm katana được đeo trong vỏ trên lưng của em. Cô nàng còn đi đôi tất cao đến đùi màu trắng với chi tiết vàng và xanh ngọc, cùng đôi dép khá giống với dép của các samurai xưa kia.

\chrint

Thôi thì, vạn sự khởi đầu nan. Chỉ hy vọng sự ``trong sáng'' của Iroha sẽ giúp tôi thở phào đôi chút\dots

\il

\begin{center}
    \uline{Sa mạc.}
\end{center}

Nakiri Ayame - cô nàng Quỷ sừng đến từ học viện xứ Quỷ, xuất thân từ một dòng tộc danh giá gắn liền với những bộ kimono truyền thống. Cô nàng vốn là một tiểu thư thích chơi khăm người khác nhưng lại cực kỳ dễ bị chọc lét, đồng thời cũng mang danh ``hội trưởng hội học sinh'' quyền uy của ngôi trường.

Đối thủ lần này của cô chính là Kazama Iroha. Một cô nàng samurai trẻ tuổi bước ra từ vùng núi sâu hoang vắng, sát cánh cùng chú chồn nhỏ Pokobe. Iroha nhận lời gia nhập tổ chức tội phạm holoX cốt chỉ để kiếm tiền mưu sinh qua ngày, ngẫm lại cũng thấy thật cơ cực làm sao. Ngoài ra, cô nàng còn nổi tiếng là một kẻ cực kỳ ``PON'', dù độ ngốc nghếch xem ra vẫn chưa chạm đến tầm\dots\ à không, cũng chẳng kém cạnh gì Miko là mấy.

\img{13-4a}

Và rồi\dots\ rốt cuộc thì tôi đang đứng ở xó xỉnh nào thế này nhỉ? Phải rồi, là bãi cát của trận đấu.

Chẳng là, ngay từ khoảnh khắc cô nàng samurai đặt chân đến văn phòng, thay vì chào hỏi lịch sự như những thành viên bình thường khác\dots

\begin{center}
    \uline{Hồi tưởng.}
\end{center}

Bên trong văn phòng, Iroha xông thẳng vào cửa với tốc độ chóng mặt, chẳng thèm chừa cho chúng tôi lấy một giây để mở lời: ``Ồ, chào buổi sáng, Iroha-cha--''

Iroha lớn giọng cắt ngang: ``Ayame-senpai! Em chính thức thách đấu chị trận đấu kiếm phân định cao thấp này, de gozaru!''

Ayame ngơ ngác chớp mắt: ``Ơ\dots\ Ể\dots\ Ể? Chị á!?''

Tôi liền hoảng hốt can thiệp: ``Khoan, khoan, khoan đã nào! Sao tự dưng chưa kịp chào hỏi câu nào mà đã xông vào thách đấu người ta thế này?''

Game thở dài qua cổ tay: ``Ít ra có chuyện gì thì cũng phải từ từ ngồi xuống uống ngụm nước đã chứ.''

Cô nàng Samurai kiên quyết lắc đầu: ``Nhưng mà máu samurai trong người em sôi sục lắm rồi, không thể chờ đợi thêm được nữa, de gozaru! Em muốn đọ thử thanh kiếm của chị ấy xem thanh nào xịn và bền bỉ hơn, gozaru!''

Ayame cạn lời: ``\dots Thật luôn hả trời?''

Tôi nhăn mặt: ``Nghe xàm xí đú thế nào ấy\dots''

Dẫu vậy, nàng Quỷ sừng máu chiến vẫn vui vẻ gật đầu nhận lời thách đấu với ``lính mới''. Trời ạ, cái văn phòng này chưa ngày nào là không nổi lên những trận hỗn chiến vô tri thế này cả\dots

\begin{center}
    \uline{Kết thúc hồi tưởng.}
\end{center}

Đó chính là lý do vì sao giờ đây chúng tôi lại đang đứng chình ình giữa một khung cảnh sa mạc hoang vu thế này. Còn vì sao lại là sa mạc á? Chắc là vì\dots\ nó mang tính nghệ thuật phù hợp với bối cảnh kiếm đạo chăng? Có giời mới biết được dụng ý nghệ thuật quái quỷ gì ở đây.

Tôi đành miễn cưỡng đứng ra cầm chiếc quạt giấy (thay vì cầm cờ) để làm trọng tài cho trận đấu định mệnh này. Mặc dù cho đến tận bây giờ, tôi vẫn chẳng tài nào hình dung nổi trận chiến này rồi sẽ đi về đâu.

E hèm\dots

Thứ duy nhất tôi có thể xác định lúc này là hai con người ấy đã trừng mắt lườm nguýt lẫn nhau suốt một khoảng thời gian dài đằng đẵng, hai tay lăm lăm chuôi kiếm với vẻ mặt cực kỳ cảnh giác. Trông bọn họ thận trọng chẳng khác nào sát thủ chuyên nghiệp. Chẳng thể đoán trước được ai sẽ ra tay trước, nhưng nếu cứ tiếp tục giằng co thế này thì đến cả tôi đứng ngoài xem cũng thấy căng thẳng sắp đứt mạch máu não.

Cộng thêm cái thời tiết nóng nực như muốn bốc hỏa giữa lòng sa mạc nữa chứ.

Game càu nhàu: ``Err\dots\ Chúng ta còn phải đứng phơi nắng thế này đến bao giờ nữa đây?''

Tôi lau mồ hôi hột: ``Chịu chết\dots\ Thề luôn, thời tiết nóng vãi chưởng\dots''

Ông ta tiếp tục phàn nàn: ``Ta đã đứng cắm rễ ở đây tròn một tiếng đồng hồ rồi đấy. Đến khi nào họ mới chịu động thủ vậy\dots?''

Ryuuhou phe phẩy tay trước mặt, vừa né nắng vừa ngó hai người trên bãi cát cằn nhằn: ``Nắng muốn khét cả da rồi Onii ơi! Em tưởng ra xem đấu kiếm đỉnh cao thế nào, ai dè hai bà chị này tính thi xem ai đứng tạo dáng bất động lâu hơn à? Nếu còn cù nhầy nữa khéo em phải dựng lều cắm trại mất.''

Bất thình lình, không để ai kịp trở tay, cả hai bóng người vút lên lao thẳng vào nhau như chớp, dường như họ đã bắt được sơ hở chết người của đối phương.

*VÚT! KENG!*

Hai bóng người lướt qua nhau nhanh như một cơn gió lốc, kèm theo âm thanh kim loại va chạm sắc lẹm xé toạc không trung. Cùng lúc đó là một tiếng cười khẩy tự tin của kẻ tân binh và tiếng nuốt nước bọt khan đắng ngắt của bậc lão luyện.

Và kết quả diễn ra sau đó khiến tôi phải trợn tròn mắt kinh ngạc, bởi vì trận chiến chớp nhoáng ấy đã kết thúc nhanh đến mức khó tin: chiếc mặt nạ quỷ đặc trưng trên đầu nàng tóc trắng bỗng bị chẻ làm đôi một cách gọn ghẽ, và hai thanh kiếm giấu sau lưng cũng đồng loạt gãy gục thành hai đoạn. Cả tôi lẫn Game đều chết lặng trước kết cục quá phũ phàng này.

\img{13-4b}

Ryuuhou đứng bên cạnh huýt sáo trêu chọc: ``Oa, một cú lội ngược dòng chớp nhoáng! Kiếm của Ayame-san giòn như bánh đa vừng vậy á?''

``Và người giành chiến thắng chung cuộc là\dots\ Kazama Iroha-chan!'' tôi vung chiếc quạt trong tay sang phía bên trái, dõng dạc tuyên bố người chiến thắng trong cuộc đối đầu chớp mắt vừa rồi.

Cô nàng samurai nhảy cẫng lên vì sung sướng, giơ cao tay hét lớn: ``Em thắng rồi, de gozaru!'', trái ngược hoàn toàn với vẻ mặt tiu nghỉu, khóc không ra nước mắt của Ayame.

Ryuuhou bước đến vỗ vai an ủi cô nàng Quỷ sừng đang buồn bực: ``Thôi Ayame-san đừng buồn nữa, dù sao thì chị cũng\dots\ đứng tạo dáng ngầu được hẳn một tiếng đồng hồ mà.''

Game cũng hùa vào góp ý tỉnh bơ: ``Giá mà cô chịu khó đem kiếm ra mài mòn sắc bén mỗi ngày thì đâu đến nỗi thua muối mặt thế này\dots''

Ayame đổ mồ hôi hột, mếu máo: ``Đó\dots\ đó không phải là vấn đề chính ở đây\dots''

Còn tôi thì nhăn mặt càu nhàu với hai người: ``Lần sau muốn góp ý thì chọn cái gì mang tính xây dựng hữu ích hơn một chút xem nào.''

\ct

Bốn người chúng tôi ngồi lại quây quần để trò chuyện giao lưu làm quen. Đặc biệt là với cô nàng Oni vốn rất hay cười tít mắt nhưng bản chất lại vô cùng hướng nội.

Mà nhắc đến chuyện cười, em ấy dường như đã vui vẻ hớn hở trở lại chỉ sau vài ba câu giao tiếp xã giao đơn giản. Đúng là một cô nàng dễ dãi và ba phải thật đấy. Mặc dù trong lòng Ayame vẫn có đôi ba lần thoáng run sợ trước sức mạnh áp đảo của đối thủ vừa hạ gục mình.

Nàng Oni ngập ngừng hỏi: ``V-Vậy\dots\ cái thanh kiếm sắc bén đó của em có tên gọi riêng nào không vậy?''

Iroha tự hào đáp: ``Dạ có chứ, nó tên là Chakimaru de gozaru.''

\img{13-4c}

Tôi nhíu mày: ``Hả? Chakimaru?''

Ryuuhou buông lời nhận xét: ``Chakimaru? Sao nghe như tên một chú cún con hay loại bánh ngọt truyền thống thế nhỉ?''

Thế nhưng, cái tên thanh kiếm của Iroha thực chất vẫn còn khá bình thường chán nếu đem so sánh với bên còn lại. Ayame không ngần ngại rút ra hai thanh katana từ sau lưng (những thanh kiếm mà tôi vừa cặm cụi ngồi sửa giúp cô ấy ban nãy) rồi kiêu hãnh tuyên bố cái tên khiến tôiém chút ngã ngửa ra đằng sau: ``Thanh kiếm bên trái của chị chính là Quỷ kiếm Asura, còn thanh bên phải là Ảo kiếm Rakshasa đấy!''

\img{13-4d}

Game thở dài bó tay thốt lên trong sự kinh ngạc tột độ: ``Cái bộ tên kiếm của cô mới là có vấn đề nặng về thần kinh thẩm mỹ đấy!''

Ryuuhou bật cười, nghiêng đầu nhìn Ayame: ``Nghe oai thì có oai thật đấy, nhưng chị vừa thua đo ván người cầm Chakimaru đấy nhé, Ayame-san.''

Thế nhưng nàng tóc vàng lại chẳng hề tỏ ra ngạc nhiên chút nào, vẫn giữ nguyên khuôn mặt ngây thơ rạng rỡ: ``Thật vậy sao ạ? Nghe oai phong lẫm liệt lắm chứ, de gozaru!''

Nói rồi, em ấy liền quay ngoắt đầu lại, hớn hở cắm cổ chạy về phía bờ biển, miệng reo lên: ``Em phải chạy đi khoe chiến tích chiến thắng này cho mọi người ở nhà biết ngay mới được!''

Đợi đã\dots\ Ờm\dots\ Lời nói đó nghĩa là em ấy có người thân đi cùng đến đây sao?

Ayame nghe thấy thế thì hoảng hồn hốt hoảng, vội vã lao vọt ra đằng trước chặn đứng đường đi của Iroha lại: ``Ấy ấy! Đứng lại ngay! Tuyệt đối không được phép nói cho ai biết chuyện này nghe chưa!''

Đây chính là lần đầu tiên trong đời tôi được chứng kiến mà Ojou rơi vào trạng thái hoảng loạn tột độ như thế. Không lẽ\dots

Iroha ngơ ngác chớp mắt: ``Tại sao lại không được nói ạ? Thắng làm vua thua làm giặc mà, de gozaru!''

Ayame đỏ bừng mặt quát: ``Làm thế là mất hết thể diện danh dự của đại tiểu thư nhà này trước mặt đàn em đấy hiểu chưa!''

Ryuuhou chống nạnh cười khì khì: ``Ayame-san tính ra cũng sĩ diện gớm.''

Game cũng chẳng kém cạnh mà chêm vào: ``Thua cuộc thì cứ dõng dạc nhận thua thôi chứ có gì mà phải giấu giếm. Chuyện bại tử binh gia là bình thường mà.''

Tôi thở dài: ``Em làm như bản thân thực sự mang cốt cách samurai chân chính không bằng ấy\dots''

Ayame xua tay liên tục: ``Không, không, không phải thế. Mọi người chẳng hiểu gì sất.''

Tôi dám chắc là Ayame đang cố vớt vát chút liêm sỉ cuối cùng của mình, nên mới phải ngăn cản quyết liệt đến vậy. Thôi thì, cứ để cho em ấy tự vẫy vùng giải thích trước đã, dù lúc đó Iroha đang tỏ thái độ không phục ra mặt. Ojou hắng giọng, cố ra vẻ bề trên: ``Cứ bình tĩnh và để đại tỷ ta đây chỉ dạy cho điều này.''

Ayame nghiêm túc ra giá: ``Giữ kín bí mật thể diện này cho ta, rồi sau này ngươi muốn đòi đền đáp xứng đáng thế nào ta cũng chiều.''

Tôi châm chọc: ``\dots Nghe điều kiện kiểu gì cũng thấy không đáng tin cậy tí nào cả.''

Game lạnh lùng: ``Hơn nữa tôi không nghĩ Kazama-san sẽ chịu nghe lời dụ dỗ đâu. Cô đang đàm phán ở tư thế của kẻ thua cuộc đấy Ojou ạ.''

Quả nhiên đúng như dự đoán. Ngay khi nàng Quỷ sừng vừa dứt lời dụ dỗ, một tiếng còi tàu hơi nước trầm đục bỗng vang lên từ phía xa xôi thu hút sự chú ý của tất cả chúng tôi.

Iroha reo lên: ``Ể? Em nghe quen quen cái tiếng còi tàu này lắm nhé!''

Ngay trước mắt, từ ngoài khơi xa, một người đàn ông mặc quân phục có khuôn mặt lai giống người châu Âu đang đứng hiên ngang trên một chiếc thuyền buồm gỗ cổ kính được đóng từ tận thế kỷ XV, XVI nào đó. Chiếc thuyền buồm kỳ lạ ấy cứ thế trôi nổi phập phồng ngay trên vùng nước bờ biển thời hiện đại và lướt qua mặt chúng tôi.

\img{13-4e}

Nàng Samurai hớn hở chỉ tay: ``Nhìn kìa! Là thuyền của Đề đốc Perry kìa, de gozaru!''

Ayame và tôi đồng thanh hét lên hụt hẫng: ``Perry cơ á!?''

Game cất giọng thắc mắc qua cổ tay: ``Ý cô là Matthew Calbraith Perry, vị phó đề đốc hải quân lừng danh lịch sử đó á hả?''

Ryuuhou tròn xoe mắt, dụi dụi mi tâm như không tin vào những gì đang thấy: ``Ủa khoan\dots\ em có bị say nắng nhìn gà hóa cuốc không vậy Onii? Đề đốc Perry lái thuyền chiến thời Edo tới bãi biển thế kỷ hai mươi mốt để hóng drama thật đấy à!?''

Iroha hào hứng: ``Đúng thế, con Hắc Thuyền huyền thoại đã cập bến rồi kìa de gozaru!''

Ayame vã mồ hôi: ``Ông ta đời thuở nào làm gì có sống chung thời kỳ canh tân với chúng ta đâu hả!''

Tôi cạn lời: ``Lại còn dạt vòm bén mảng lái thuyền chiến tới tận bờ biển thế kỷ XXI này lướt sóng nữa chứ, thật là loạn lạc!''

Người đàn ông lạ mặt đứng trên chiến thuyền cất giọng hỏi vọng vào bờ: ``Này mấy đứa nhóc bờ biển, kẻ nào vừa giành chiến thắng trong trận đấu kiếm thế?''

Tôi vô thức đổ mồ hôi hột, thầm nghĩ trong đầu: ``\textit{Đến cả lão hải quân ngoại quốc này mà cũng hóng biến tấu hài nữa cơ à?}''

Iroha hí hửng giơ tay hét lớn đáp lại: ``Là tôi chứ ai vào đây n--''

Ngay lập tức, em ấy bị đôi tay nhanh hơn não của Ayame chồm tới bịp chặt cứng miệng lại. Nàng tóc trắng trợn tròn mắt tỏ vẻ hăm dọa đầy hoảng sợ: ``Đừng nhìn điều xấu, đừng nghe điều xấu, và tuyệt đối không được `de gozaru' hé răng nhắc đến kết quả xấu hổ đó cho ông ta nghe rõ chưa hả! Nè!?''

\img{13-4f}

Nàng tóc vàng ngơ ngác đến mức quá đỗi ngỡ ngàng trước phản ứng bộc phát thái quá của đàn chị, đến mức ngớ người không biết phải đáp lại thế nào: ``Hể\dots?''

Ryuuhou đứng bên cạnh ôm bụng cười khúc khích: ``Nhìn kìa, Ayame-san cuống cuồng bịt miệng đàn em đến toát cả mồ hôi hột rồi kìa! Sợ mất mặt với cả danh nhân lịch sử cơ đấy!''

Còn tôi và Game thì chỉ biết thở dài ngao ngán: ``Có cái quái gì mà phải sợ mất mặt thế cơ chứ\dots'' / ``Sớm muộn gì thì cái thiên hạ này người ta cũng tự khắc biết hết thôi.''

Tất nhiên là Ayame sau đó cũng quay sang trừng mắt hăm dọa cả đám chúng tôi, nhưng trông bộ dạng lúc ấy thảm thiết và tội nghiệp lắm rồi: ``Xin hai người đấy, coi như em lạy, tuyệt đối đừng buột miệng nói cho bất kỳ ai biết về chuyện này nhé, yo!''

Ojou-san đã hạ mình cầu xin đến mức độ đó thì tôi đành phải ngậm bồ hòn làm ngọt đồng ý vậy.

Chậc. Đúng là cái thứ phiền phức không bao giờ để cho người ta được bình yên mà.

\ct

Chúng tôi lết xác quay trở lại văn phòng, tận hưởng luồng gió điều hòa mát lạnh phả ra sau khi phải đứng phơi nắng và đi bộ rã rời suốt hai tiếng đồng hồ ngoài bãi cát. Cuối cùng thì tôi cũng có thể thả lỏng cơ thể, trở lại với trạng thái lười biếng tự nhiên thường ngày của mình.

Ryuuhou vừa ngồi phịch xuống ghế vừa thở dài sườn sượt: ``Phù\dots\ sống lại rồi! Lần sau Onii có đi làm trọng tài giữa sa mạc thì tự thân vận động nhé, em không dại gì đi phơi nắng nướng đen cả người đâu.''

Hôm nay không có quá nhiều việc gấp, thế nên tôi đang ngồi thõng người ở một đầu của chiếc ghế bành phòng khách, cặm cụi lướt điện thoại, trong khi Iroha thì đang ngồi nghiêm túc ở đầu ghế bên kia, hì hục lật danh sách các tài năng của công ty để tìm kiếm một đối thủ xứng tầm tiếp theo.

Chẳng hiểu tại sao con bé samurai này lại cứ thích đi khiêu chiến khắp nơi như thế không biết.

Iroha ngẩng đầu lên, lí nhí gọi: ``Mà\dots\ Kuon-dono\dots''

Tôi hờ hững đáp: ``Hửm? Gì thế?''

Iroha chớp mắt hỏi: ``Anh nghĩ trong số tất cả mọi người ở đây, ai mới là kẻ đáng gờm và ghê gớm nhất ạ?''

Game cất giọng thắc mắc: ``Cô đang định đánh giá dựa trên khía cạnh bạo lực hay trí tuệ thế? Cơ mà, từ khi nào mà cô lại bắt đầu gọi Kuon-user là `Kuon-dono' thế hở?''

Đó cũng chính là cách mà con bé xưng hô với tôi, hệt như cách mà một bề tôi trung thành gọi lãnh chúa của mình vậy. Tôi thậm chí còn chẳng phải là một nhân vật lớn lao gì sất để mà nhận cái danh xưng trang trọng kiểu đó.

Tôi nghĩ chị A-san hay chú Yagoo (Tanigo-san) mới là những người thực sự xứng đáng với cái danh hiệu ``-dono'' tôn nghiêm ấy hơn. Thế nhưng ngặt nỗi, em ấy cứ khăng khăng muốn gọi tôi như thế chỉ vì tôi là Tổng quản lý duy nhất của cả nhánh \hll\ này. Một lý do nghe chẳng có chút sức thuyết phục nào cả.

Ngay lúc đó, Ryuuhou từ góc phòng ngẩng đầu lên, tặc lưỡi đệm vào: ``Gọi là `ông anh lười' hay `quản lý quèn' cho nó thực tế Iroha-san ơi, chứ `Kuon-dono' nghe cứ như lãnh chúa thời Chiến quốc ấy, Onii nhà em làm gì có cửa oai phong lẫm liệt thế!''

Đúng lúc ấy, một giọng nói dè dặt vang lên từ giữa phòng: ``Ờm\dots\ hai người cứ coi như không nhìn thấy sự tồn tại của em ở đây nha\dots\ Em chỉ nấp ở góc này thôi\dots''

Tôi liếc mắt sang thì thấy Ayame đang lén lút thò đầu ra từ sau bức tường, nơm nớp nhìn chằm chằm về phía chúng tôi.

\img{13-4g}

Game gằn giọng mỉa mai: ``Cô đã rình mò bám đuôi chúng tôi suốt từ lúc bước vào văn phòng cho đến tận bây giờ rồi đấy, Nakiri-san ạ. Cô thật sự rảnh rỗi và dai như đỉa đó thế cơ à?''

Hiển nhiên là Iroha cảm thấy cực kỳ không thoải mái khi cứ bị theo dõi như thế. Cô nàng liếc mắt sang với vẻ mặt bất mãn: ``Chị làm ơn đừng có đi rình mò bám đuôi bọn em nữa chứ!''

Tôi thở dài thanh minh: ``Anh đã hứa chắc nịch với em là sẽ giữ kín bí mật, tuyệt đối không hé răng nói với ai về chuyện trận thua muối mặt ngoài bãi cát rồi cơ mà.''

Nàng Oni mếu máo bám víu: ``Nhưng đấy là anh hứa! Nhỡ đâu con bé lỡ miệng đi khoe với người khác thì sao!?''

Iroha giơ hai tay đan chéo nhau lại thành hình dấu X trước ngực, kiên quyết khẳng định rằng bản thân vẫn đang giữ đúng lời hứa danh dự với tiền bối: ``Em đã hứa là giữ kín rồi cơ mà! Không thấy gì, không nghe gì, không nói gì cả!''

\img{13-4h}

Nhưng đáp lại sự quả quyết ấy chỉ là cái nghiêng đầu ngơ ngác khó hiểu từ Ayame: ``Cái quái gì thế\dots?''

Cả tôi và Iroha bất lực kêu lên: ``Chính em/chị vừa oang oang cái mồm đi nhắc lại chuyện đó với em (ấy) chứ còn ai vào đây nữa!''

Game bồi thêm: ``Tự mình bóp chát chính mình mà không hề hay biết luôn kìa.''

Ryuuhou nhếch mép cười trêu: ``Chuẩn luôn, Ayame-san tự đào hố chôn mình sâu cả thước rồi kìa!''

Bất thình lình, Haato từ xó nào lao vù ra như một con thú hoang, nhảy bổ vào đầu nàng samurai: ``Ồ, hố hố!'' - cô nàng tóm chặt lấy mái tóc của Iroha, ghì chặt không buông.

\img{13-4i}

Cả nàng Quỷ sừng lẫn chúng tôi đều trố mắt nghệt mặt ra, hoàn toàn đơ người không hiểu chuyện quái gì đang xảy ra trước mắt.

Game cạn lời: ``Cô đang giở cái trò tâm thần phân liệt gì thế hả, Akai-san\dots?''

Iroha mặt xanh như tàu lá, hoảng hốt kêu la: ``\textbf{Ái da, đau, đau quá! Kuon-dono, Ayame-senpai, cứu em với!}''

Ayame ngơ ngác chớp mắt: ``Ể? Chuyện gì thế này?''

Tôi nhíu mày can thiệp: ``\dots Nghiêm túc đấy Haato-chan, bộ em đói khát va chạm lắm rồi hay sao mà cứ sồn sồn lao vào bạo hành lính mới của chúng ta thế hả?''

Phải mất một lúc vật lộn, tôi mới kéo được Haato ra khỏi đầu nạn nhân. Nàng tóc vàng Gen 1 hắng giọng ra vẻ bề trên rồi dõng dạc tuyên bố: ``Chị nghe ngóng được tin đồn là em đang lùng sục đi tìm một đối thủ thực sự đáng gờm đúng không?''

Ngay lập tức, Iroha vẫn với khuôn mặt nhăn nhó vì đau, vội lắc đầu nguầy nguậy: ``Cái sự đáng gờm của chị chắc chắn là khác hoàn toàn với tiêu chuẩn bình thường rồi!''

Ryuuhou hùa theo bồi một đao: ``Người Nhện phiên bản lỗi thì có, nếu không muốn nói là dị hợm không ngửi nổi luôn á Haato-san.''

Như muốn chứng minh năng lực của bản thân, toàn thân Haato bỗng tỏa ra một thứ ánh sáng màu hồng kỳ dị. Cô nàng bắt đầu khua tay múa chân thực hiện mấy động tác ``vận khí công'' lố lăng.

\img{13-4j}

Haato ngẩng cao đầu hô lớn: ``Chị đây chính là một trong những bậc thầy sử dụng nhện điêu luyện và tàn bạo nhất ở xứ sở Phù Tang này đấy! Mọi người trong giới giang hồ thường nể sợ gọi chị là Người Nhệ--''

Game phũ phàng cắt ngang: ``Cô đang lảm nhảm cái quái gì thế?''

Tôi châm chọc: ``Ai mà chẳng lạ gì em, ngoài cái mồm điêu ngoa bốc phét ra thì có cái tích sự gì đâu.''

Thậm chí, Iroha còn hoảng hốt đưa hai tay lên bịt chặt tai lại, hét toáng lên đầy sợ hãi: ``\textbf{Em không nghe, em không biết, em không thấy gì hết!}'' rồi cắm cổ chạy một mạch trối chết khỏi văn phòng, tốc độ nhanh đến mức chúng tôi chẳng kịp trở tay ngăn cản.

\img{13-4k}

Tôi giật mình kêu lớn: ``Ơ kìa, này Iroha! Chạy đi đâu thế hả!?''

Haato phất tay chảnh chọe: ``Đúng là một đứa trẻ thiếu lễ phép thật đấy. Cả gan dám chen ngang lời của tiền bối cơ đấy\dots''

Ayame liếc xéo, mỉm cười khinh bỉ: ``Em nghĩ là chị còn vô lễ và mất nết hơn khi dám lao vào cắn xé đầu tóc của con bé đấy.''

Game đồng tình: ``Hoàn toàn chính xác. Đứng ngẩn ra làm gì nữa, đuổi theo cô ta nhanh lên thôi, Kuon-user.''

Ryuuhou bật dậy vươn vai hào hứng: ``Hóng biến tiếp thôi! Để xem bạn samurai trốn đi đâu được nào!'' Con bé gọi với theo, giọng đầy vẻ trêu chọc: ``Iroha-san, chạy chậm thôi! Chờ bọn này với, để xem chị định thách đấu ai tiếp theo nào!''

Tôi thở dài ngao ngán: ``Phải chịu thôi\dots\ Này, chờ bọn anh với, Iroha-chan!''

\ct

Iroha cứ thế cắm cổ vừa chạy vừa bịt chặt hai tai, chạy một mạch ra khá xa khỏi khu vực văn phòng. Thế nhưng, trông em ấy vẫn phơi phới, dường như chẳng hề thấy mệt mỏi chút nào. Quả không hổ danh là một samurai được rèn luyện thể lực từ chốn thâm sơn cùng cốc có khác.

Còn tôi thì đã bắt đầu thở phì phò, lè lưỡi như một con bò tót sau màn đuổi bắt vừa rồi. Ryuuhou thong thả đi sát bên cạnh, vừa tung hứng một hòn sỏi nhỏ vừa châm chọc: ``Mới chạy có một đoạn mà Onii đã thở dốc như kéo bễ rồi kìa, thể lực yếu thế này bảo sao suốt ngày bị mấy bà chị ở công ty đè đầu cưỡi cổ.''

Iroha vừa chạy vừa quay đầu lại hỏi: ``Chúng ta\dots\ đã cắt đuôi được chị ta chưa vậy Kuon-dono?''

Tôi chống gối thở hồng hộc: ``Còn\dots\ còn hơn cả cắt đuôi nữa chứ. Em chạy kiểu gì mà làm anh mệt đứt cả hơi thế này\dots''

Nàng Samurai hối lỗi: ``Em xin lỗi nha, Kuon-dono!''

Tôi xua tay: ``Bình thường thôi em. Anh toàn bị mấy đứa trong công ty hành hạ thế này suốt quen rồi. Mà\dots\ làm ơn đổi cái cách gọi `Kuon-dono' ấy đi được không, nghe cứ như phim cổ trang lịch sử ấy\dots''

Ryuuhou cười hì hì: ``Thôi chấp nhận đi Onii, người ta đã tôn anh làm chủ tướng rồi đấy, hôm nào mua bộ áo giáp samurai về mặc cho hợp vai nhé!''

Bất chợt, một giọng nói trong trẻo bỗng vang lên từ phía trên đỉnh đầu: ``Iroha-chan!''

Chúng tôi giật mình ngước nhìn lên phía phát ra âm thanh.

Iroha ngạc nhiên: ``Roboco-senpai?''

Tôi nhíu mày: ``Em đang làm gì mà lơ lửng ở trên đấy thế, Roboco?''

Roboco từ trên cao từ tốn đáp xuống, nhún vai: ``Em nghe ngóng được tin là có người đang đi lùng sục khắp nơi để khiêu chiến. Iroha-chan muốn tìm đối thủ mạnh để đấu kiếm, đúng không nào?''

Tôi thở dài: ``Thì đúng là vậy thật, nhưng mà\dots''

Iroha hừng hực khí thế, giơ hai nắm đấm lên cao: ``Vậy chị định dùng chiêu thức gì để chiến đấu với em đây?''

Roboco mỉm cười giơ phắt bàn tay trái lên. Ryuuhou tròn mắt nhìn bàn tay phát sáng tím lịm, hoang mang thốt lên: ``Ủa khoan\dots\ Roboco-san định đấu kiếm bằng năng lượng hạt nhân à!?''

Game cất giọng cảnh báo qua cổ tay: ``\dots Tôi có linh cảm không lành, khuyên cô tốt nhất là không nên hỏi câu đó đâu.''

Luồng năng lượng trong lòng bàn tay Roboco tụ lại, dần dần sáng rực lên đến chói mắt.

\img{13-4l}

\dots Chết tiệt! Lại là cái điệp khúc quen thuộc này nữa chứ!

Em ấy lại chuẩn bị dùng tia laser hủy diệt rồi! Đúng là cái đồ đại PON không thể cứu vãn nổi mà! Tôi hoảng hốt thét lớn: ``\textbf{Nằm xuống!! Chạy!!}''

*BÙM! CHÍU!*

Vừa dứt câu, Robo-PON đã bóp cò phóng thẳng một chùm tia laser hủy diệt từ lòng bàn tay ra, không may nã trúng ngay một thùng xăng dự trữ bên đường khiến nó phát nổ tung trời, sức ép thổi văng chúng tôi bay tút ra xa.

Iroha hoảng loạn hét lên: ``\textbf{Đòn tấn công của chị ấy hiện đại và phi lý quá mức cho phép rồi đấy, de gozaru!}''

Ryuuhou vừa nhảy lùi lại vừa xua khói bụi cay xè mắt: ``Má ơi! Kiếm đạo này lạ quá, samurai thời phong kiến nào đỡ cho thấu trời!''

Game gằn giọng phàn nàn: ``Đừng có mà vừa di chuyển vừa nã súng laser bừa bãi thế chứ, cái đồ óc bã đậu công nghệ cao kia!''

Iroha vừa chạy vừa mếu máo: ``Em đến từ một vùng núi sâu bị cách biệt hoàn toàn với thế giới văn minh Nhật Bản hiện đại cơ mà, sao đỡ nổi mấy thứ này!''

Tôi vừa né đạn vừa cáu kỉnh gào lên: ``Em làm thế thì bố con anh có chín mạng cũng chẳng đủ để đọ lại với cái độ tàn phá của em đâu!''

Cô nàng Người máy vẫn liên tục nhắm chuẩn nã laser về phía chúng tôi, tỏ ra vô cùng hăng hái không hề có ý định khoan nhượng. Bản thân tôi và Ryuuhou cũng suýt thì dính đạn oan gia vài ba lần từ những phát bắn lệch quỹ đạo của em ấy.

Thế nhưng rắc rối chưa dừng lại ở đó. Bất ngờ thay, một thành viên \hll\ khác lại ngoi lên từ dưới chiếc nắp cống bên đường, khiến chúng tôi suýt đâm đầu vào nhau phải ``phanh gấp'' trượt chân ngã dúi dụi.

Đó chính là Noel, người hiện tại đang cosplay đóng giả làm một thây ma thực thụ.

``Thức ăn\dots\ Lẹ mang thức ăn dâng lên cho ta mau!'' hoặc nói một cách chính xác hơn, đây là một con thây ma đói lả, thay vì đi săn não người sống thì lại đang gào thét đòi ăn một phần cơm bò gyudon vì đã bị bỏ đói suốt cả ngày trời.

\img{13-4m}

``Thây ma ư? Ể\dots\ Chẳng phải đó là Noel-senpai đó sao?'' Nàng samurai ngơ ngác hỏi tôi. Tôi ngán ngẩm thở dài: ``Chứ còn ai trồng khoai đất này nữa.''

Ryuuhou lùi lại một bước, chỉ tay cười trừ: ``Ái chà, zombie Noel-san trồi lên từ cống ngầm kìa! Nhìn cái bộ dạng đói khát này chắc chắn là chưa được ăn trưa rồi! Iroha-san chạy mau không là bị đem ra làm thức ăn đấy!''

Nàng Hiệp sĩ thây ma cất giọng khàn đục như người chết đói (dù nghe qua vẫn thấy độ giả trạo không che giấu nổi): ``Ta ngửi thấy\dots\ mùi vị của một cô nàng samurai thơm ngon béo bở đâu đây\dots''

Dường như bản năng ăn uống kết hợp với máu chiến đã khiến em ấy lập tức tia trúng Iroha làm con mồi.

Noel bỗng nhiên phát điên lên thực sự, gào thét rống lên như quái vật rồi điên cuồng lao thẳng về phía chúng tôi. Thật sự luôn, rốt cuộc là kẻ tàn nhẫn nào đã bỏ đói em ấy suốt mấy ngày trời thế hả trời?

Nhưng dù lý do là gì đi nữa, con đường này coi như đã bị chặn đứng hoàn toàn!

Tôi hoảng loạn quay đầu tìm lối thoát, bất chợt nhìn thấy nàng Người máy Roboco ở phía sau đang đứng đực ra vẻ mặt hốt hoảng. Có vẻ như em ấy đang lo lắng cho sự an toàn của Noel hơn là bận tâm đến việc truy sát cô nàng tóc vàng.

Iroha hét toáng lên: ``\textbf{Chị ấy đang lao thẳng về phía chúng ta kìa, de gozaru!}''

Tôi dứt khoát ra lệnh: ``Bám chặt lấy áo anh! Để anh lo việc này! Game, kích hoạt lớp khiên năng lượng tối đa ngay lập tức!''

Game đáp gọn: ``Rõ, Kuon-user!''

Thế nhưng, một giọng nói uy nghiêm khác cất lên đúng lúc đã cắt đứt hoàn toàn cuộc hỗn chiến vô tri ấy. Chúng tôi giật mình quay phắt đầu ra đằng sau, và đập vào mắt là hình ảnh Sora đang đứng chống nông, nghiêm nghị mắng mỏ hai cô nàng đang cúi gằm mặt vì sợ hãi.

Sora cất giọng lạnh lùng: ``Đừng có mà hù dọa và đi săn lùng thành viên mới của chúng ta như thế chứ, hư quá đấy.''

\img{13-4n}

Noel và Roboco đồng thanh rối rít xin lỗi: ``Bọn em/tôi xin lỗi Sora-senpai/Sora-chan\dots''

Ryuuhou khoanh tay tán thưởng đầy thán phục: ``Ối chà, đúng là Sora-san có khác! Vừa xuất hiện một cái là hai chiến thần dập tắt ngấm lửa luôn, uy quyền ngút trời!''

Tôi chép miệng lẩm bẩm: ``Đúng là `vỏ quýt dày có móng tay nhọn', một đòn uy quyền ăn ngay lập tức kìa\dots''

Game gằn giọng: ``Quả không hổ danh là thần tượng đời đầu tối cao của công ty\dots''

Nàng Thần tượng hoàn hảo từ từ quay mặt về phía chúng tôi, nhẹ nhàng chắp tay cúi đầu xin lỗi, xung quanh cơ thể toát ra một thứ ánh hào quang thanh tịnh đến lạ thường.

Sora dịu dàng nói: ``Xin lỗi mọi người vì rắc rối vừa rồi nhé.''

Iroha ngượng ngùng đáp: ``Dạ\dots\ không có gì phiền đâu chị ạ. Bọn em phải cảm ơn chị vì đã ra tay giải cứu mới đúng ạ.''

Tôi thở phào: ``Còn hai đứa chiến thần PON kia thì cứ để anh với Game lo thu xếp.''

Noel và Roboco cãi khất: ``Bọn em không có PON nhé!''

Game chê bai: ``Yeah, yeah, sao cũng được, cãi cố làm gì.''

Iroha và Sora mỉm cười nhẹ nhàng với nhau, trao nhau những ánh nhìn thân thiết ấm áp nhất. Nàng samurai thậm chí còn đỏ bừng cả hai má vì quá e thẹn trước đàn chị.

Nàng tóc nâu tò mò lên tiếng hỏi: ``Cơ mà, tại sao mọi chuyện lại trở nên rắc rối thế này vậy?''

Lập tức, máu hiếu chiến của cô hậu bối trỗi dậy, đôi mắt sáng lấp lánh hân hoan đáp ngay: ``Dạ báo cáo, em đang thực hiện một nhiệm vụ tối quan trọng là đi tìm và đánh bại kẻ mạnh nhất ở cái công ty \hll\ này ạ\dots''

Ryuuhou trố mắt, ghé sát tai tôi thì thào: ``Gan con bé này to bằng trời rồi Onii ơi, dám dõng dạc tuyên bố câu đó ngay trước mặt Boss tối cao luôn kìa\dots''

Tôi và Game cùng đổ mồ hôi hột, thầm nghĩ: ``\textit{Để làm cái quái gì cơ chứ!?}''

Nếu như con bé này mà dại dột nổi điên đi tìm chú Yagoo (Tanigo-san) để khiêu chiến, tôi dám chắc là công ty sập tiệm trước khi kịp giao đấu mất.

Thế nhưng xui rủi thế nào, vừa nghe dứt câu mục tiêu của Iroha, toàn thân Sora bỗng nhiên tỏa ra một thứ ánh sáng ma mị kỳ quái. Vẫn giữ nguyên nụ cười hiền từ (?) trên môi, Sora cất giọng ngọt ngào: ``Vậy có nghĩa là\dots\ em coi chị chính là kẻ thù cần phải vượt qua ha.''

\img{13-4o}

Tôi thừa biết ý trong câu nói của em ấy không phải là kiểu ``đánh nhau đổ máu'', mà chỉ là ``cạnh tranh lành mạnh'' trong công việc thường ngày. Thế nhưng, ngặt một nỗi là vì tiếng tăm lừng lẫy với biệt danh ``sát thủ thầm lặng'' của em ấy quá mức\dots\ đáng sợ trong công ty, nên Iroha nghe xong liền mặt cắt không còn một giọt máu, cứng họng câm nín hẳn như hến.

Game nhận xét: ``\dots Tôi nghĩ Tokino-san hoàn toàn không có ác ý gì với cô đâu. Mà này Tokino-san, cô cũng đang là người trực tiếp dọa Kazama-san đến mức hồn vía lên mây rồi kìa.''

Tôi vội can ngăn: ``Sora-chan, em làm ơn thu lại cái sát khí đó đi, đừng có dọa con bé sợ chết khiếp thế chứ\dots''

Ryuuhou ôm đầu hoảng hốt: ``Ủa ủa? Sora-san ơi, nụ cười thánh thiện của chị vừa kích hoạt hào quang làm Iroha-san hồn bay phách lạc luôn rồi kìa!''

Nghe chúng tôi nhắc nhở, nàng Thần tượng mới giật mình hốt hoảng: ``T-Thế ư!? Chị lỡ dọa em ấy sợ lắm sao!?''

Em ấy lật đật quay phắt người lại định vươn tay vỗ về Iroha để xin lỗi, nhưng ngặt nỗi nàng Samurai trợn ngược mắt, hồn vía bay sạch lên mây: ``Đáng sợ quá đi mất\dots\ Kuon-dono cứu em\dots''

\img{13-4p}

Sora hoảng hốt: ``Ể!? Iroha-chan!? Tỉnh lại đi em ơi! Đừng có dọa chị!''

Tôi lay lay người con bé: ``Tỉnh táo lại xem nào! Này!''

Game lạnh lùng chốt hạ: ``Trời ạ, cô ta thăng thiên hóa đá thực sự luôn kìa\dots''

\dots Mới vừa bước chân vào công ty chưa được bao lâu mà đã ngất xỉu ngã gục thế này. Thế mà trước đó còn lớn lối đòi đi khiêu chiến với kẻ mạnh nhất trong \hl\ cơ đấy!

Nhưng mà, suy cho cùng thì lý do thực sự đằng sau việc con bé khăng khăng đòi đi thách đấu khắp nơi là gì nhỉ?

Tôi vẫn chẳng tài nào suy luận ra nổi. Chắc chắn là em ấy đang che giấu một toan tính thâm sâu nào đó, hoặc là\dots

\dots có một thế lực tăm tối nào đó đứng sau giật dây bắt buộc em ấy phải làm thế chăng?

%==========================================TRUYỆN PHỤ=========================================
\omk

Hóa ra mọi chuyện khuất tất đằng sau đều có nguyên nhân của nó. Hoặc ít nhất, những gì tôi tận mắt nhìn thấy trong đoạn lịch sử tin nhắn giữa em ấy và Sếp của em đã sáng tỏ tất cả.

Iroha nhắn tin với một khuôn mặt hoàn toàn tiều tụy, chán chường, chẳng còn chút sức sống vui tươi giống như hồi đầu mới bước chân vào công ty.

\img{13-4q}

Về cơ bản, đoạn hội thoại chỉ đạo của bọn chúng diễn ra thế này:

\begin{center}
    \uline{Bên trái: ``Sếp'', bên phải: Iroha}
\end{center}

\tl{``Ngươi dù gì cũng là tên bảo vệ của chúng ta, thế nên ngươi thử tìm xem liệu ngươi có thể triệt hạ được đối thủ nào trong cái công ty này không.''}

\tl{``À, nhưng mà nhớ chừa cái tên Kuon kia ra rìa, không được động vào!''}

\tr{``Tại sao lại phải né anh ấy ra vậy ạ, de gozaru?''}

\tl{``Thì\dots\ cứ để đề phòng rủi ro thôi, hỏi lằng nhằng làm gì!''}

\tc{(Ngày 26 tháng 6, năm 2022)}

\tr{``Cây đao ý chí và niềm tin trong tim của mình đã bị bẻ làm đôi rồi, de gozaru\dots''}

\tl{``Hả? Ngươi đang lảm nhảm cái quái gì thế?''}

\dots Được rồi, nói thật là tôi đọc xong cũng chẳng tài nào hiểu nổi cái logic của bọn chúng. Nhưng ít nhất tôi cũng nắm được mấu chốt là do cô Sếp nhỏ bên phe holoX ép buộc em ấy phải làm vậy.

Chỉ là\dots\ tại sao cô ta lại đặc biệt dặn dò phải giữ tôi lại làm ngoại lệ nhỉ? Mục đích của con bé sừng mọc rễ đấy rốt cuộc để làm gì cơ chứ?

``Con ranh con đó trông có vẻ khả nghi và mưu mô đấy,'' Trợ lý Game gầm gừ cất giọng qua cổ tay.

Tôi chỉ biết ngửa cổ thở dài đầy bất lực, bởi vì ngoài cái tên lù tịt và bộ sơ yếu lý lịch màu mè của con bé ra, tôi chẳng có lấy một chút thông tin tình báo nào khác về Laplus-chan cả.

Dã tâm thực sự của con bé đó là gì? Có giời mới biết được!

\dots Quay trở lại với hiện thực, Iroha - người sau khi nhận được câu trả lời phũ phàng từ đối thủ định mệnh - đang chán nản ngồi ôm đầu gục xuống, toàn thân run lên cầm cập vì sợ hãi: ``Sợ quá\dots\ Sợ cái thế giới công ty này quá đi mất, de gozaru\dots''

Ryuuhou đứng bên cạnh chống hông, bĩu môi châm chọc: ``Mới va vấp có chút xíu đời mà đã co rúm lại như chú cún ướt mưa thế kia rồi, thế này thì làm vệ sĩ bảo kê cho tổ chức tà ác kiểu gì đây Iroha-san ơi.''

Chúng tôi đứng xung quanh nhìn em ấy mà không khỏi tỏ ra ái ngại, lo lắng. Có vẻ như mấy trò đùa cợt quá đà của Roboco, Noel và cả Haato vừa rồi đã đánh sập hoàn toàn tâm lý của con bé rồi.

Sora áy náy lên tiếng: ``Em thật sự không nghĩ là con bé lại dễ suy sụp và ám ảnh đến mức độ này đấy\dots''

Game bình tĩnh nhận xét: ``Có vẻ là như thế rồi. Này Kuon-user, cậu là Tổng quản lý ở đây, ra mặt làm gì đi chứ?''

Tôi nhún vai bất lực: ``Bây giờ ông muốn tôi phải đóng vai trò gì cứu vớt tâm hồn em ấy đây?''

Cuối cùng, sau một hồi đùn đẩy trách nhiệm, chúng tôi chọn giải pháp an toàn nhất là để cho em ấy tự ngồi tĩnh tâm một mình. Tạ ơn trời đất là sức sống của con bé này hồi phục tinh thần nhanh đến kỳ lạ, đúng chất một nữ samurai kiên cường.

Những kẻ vừa mới hùa vào bắt nạt, hù dọa em ấy ban nãy, thậm chí là cả tôi và Sora-chan, cuối cùng đều phải đích thân ra mặt cúi đầu xin lỗi vì đã vô tình làm ``hoen ố danh dự'' thượng võ của Kazama Iroha-chan.

Cô nàng samurai vội vã xua tay, nở một nụ cười trừ gượng gạo: ``Không, không sao đâu anh chị ạ, lỗi lơ đễnh cũng là một phần do em tự chuốc lấy thôi\dots''

Tóm lại, sau màn xin lỗi qua loa thì hai bên cũng đã giảng hòa êm đẹp. Tôi tháo chiếc smartwatch đặt lên bàn sạc rồi tạm thời ra ngoài khu vực lễ tân lấy thêm chút tài liệu họp, để lại mấy cô gái quây quần ở phòng khách. Tôi chỉ hi vọng rằng từ nay Kazama Iroha-chan có thể thích nghi và hòa nhập thật tốt với tất cả mọi người ở đây. Đó cũng là nguyện vọng chung của ban quản lý dành cho các thành viên mới.

\ct

Lúc này tại phòng khách, các cô gái đang quây quần lại với nhau để trò chuyện rôm rả về sự cố dở khóc dở cười vừa rồi.

Sora dịu dàng hỏi: ``Mà này, Iroha-chan?''

Iroha chớp mắt: ``Dạ có em?''

Haato bĩu môi chêm vào: ``Mà công nhận lúc nãy Sora-chan ra vẻ nghiêm nghị trông đáng sợ thật đấy\dots''

Nàng Thần tượng phồng má hờn dỗi: ``Này, mọi người thật sự nghĩ chị là một kẻ máu lạnh đáng sợ thế cơ à?''

Nàng Trung học tỉnh bơ nhận xét: ``\dots nhưng ngặt một nỗi là chị ấy vẫn chưa phải là người đứng đầu mạnh nhất trong cái công ty này đâu nha.''

Nàng Samurai nghe thế, tò mò nghiêng đầu: ``Vậy\dots\ rốt cuộc người mạnh nhất đó là ai ạ? Chẳng lẽ là YAGOO-san sao?''

Ayame lắc đầu nguầy nguậy: ``Không. Lão tổng giám đốc đó thì tính làm gì.''

``Thế thì\dots\ là chị A-san quyền lực ạ?''

Haato gật gù: ``Cũng xem như là gần đúng tiêu chuẩn rồi đấy\dots\ Đúng là chị ấy có nhiều pha khạc ra lửa khiến cho holomem bọn chị sợ, nhưng vẫn chưa phải.''

Sora mỉm cười thần bí: ``Có lẽ em sẽ cảm thấy bất ngờ đến rớt cằm khi bọn chị nói ra cái tên này, nhưng mà\dots''

Ayame nghiêm túc hạ giọng: ``Kuon-senpai mới chính là nhân vật đáng sợ, quyền lực và máu mặt nhất ở cả cái công ty \hll\ này đấy.''

Iroha trợn tròn mắt, há hốc mồm kinh ngạc: ``Ể\dots\ Ể!? Thật á!? Tại sao lại là anh quản lý của chúng ta chứ!?''

Ayame nhún vai: ``\dots Cô em thực sự còn non và xanh lắm, tốt nhất là không nên muốn biết lý do sâu xa đâu.''

Haato khoanh tay cảnh cáo: ``Chị khuyên thật lòng là em tốt nhất đừng bao giờ dại dột bén mảng đến làm phiền anh ấy trong lúc anh ấy đang cắm mặt làm việc thì hơn đấy\dots''

Sora gật đầu tán đồng: ``Thật đấy. Cơn giận của anh ấy mà đã bùng nổ lên rồi thì chẳng ai ở cái văn phòng này cản nổi luôn.''

Từ trên chiếc đồng hồ đeo tay của tôi đặt trên bàn sạc, màn hình bỗng lóe sáng. Trợ lý Game cất giọng trầm thấp, đều đều như một phát thanh viên đang đọc hồ sơ tuyệt mật: ``Để tôi bổ sung thêm một chút tư liệu thực tế cho cô hiểu rõ hơn về người đàn ông đó nhé, Kazama-san.''

Iroha giật bắn mình, nhìn chằm chằm vào mặt đồng hồ: ``H-Hồ sơ thực tế gì cơ, de gozaru!?''

Vị trợ lý bắt đầu tuôn một tràng với chất giọng rùng rợn, đậm mùi hù dọa: ``Kuon-user bề ngoài có vẻ chỉ là một tên quản lý lười biếng, nhưng thực chất cậu ta là kẻ đã từng sống sót trở về sau vô số chuyến chu du qua các chiều không gian dị thường.''

Iroha nuốt nước bọt cái ực: ``D-Du hành đa vũ trụ á!?''

Ông ta tiếp tục vẽ nên chân dung `bất khả chiến bại' của tôi bằng một giọng điệu cực kỳ nghiêm túc: ``Đúng vậy. Cậu ta từng bị hút vào Hư Vô (The Void) - một khoảng không vô tận do vết nứt ma thuật không gian sinh ra - mà vẫn lành lặn tìm lại và giải cứu được hai người em gái lưu lạc. Cậu ta từng xé toạc ranh giới số hóa của thế giới \textit{There Is No Game}, phá vỡ các quy tắc lập trình để lôi một trí tuệ nhân tạo như tôi ra thế giới thực này. Chưa hết, cậu ta còn từng sống sót qua mười hai ngày dài ở thế giới Cầu Vồng xa xôi\dots''

Ryuuhou khoanh tay, đắc ý gật gù chêm vào: ``Chuẩn luôn! Chính Onii đã kể như thế đấy! Mà đấy là ông ấy còn chưa kể đến vụ kinh dị nhất đâu nhé!''

Cô nàng Samurai run rẩy cả người, hai mắt chớp lia lịa: ``V-Vụ kinh dị nhất ạ\dots?''

Game tiếp lời với tông giọng lạnh tanh: ``Đỉnh điểm là tại thị trấn ma quái Aogami trong cơn ác mộng \textit{hololive ERROR}. Cậu ta bị tước đoạt toàn bộ ký ức, đối mặt với vòng lặp tử thần vô tận do một bóng ma mang diện mạo hệt như Tokino-san thao túng, vậy mà vẫn phá vỡ được kết giới nguyền rủa, hoàn thành nghi thức thanh tẩy giải thoát cho linh hồn tổ tiên cả ngàn năm tuổi đấy.''

Sora che miệng mỉm cười thần bí đầy ma mị: ``Fufu, nhắc lại chuyện ở thị trấn đó\dots\ quả là một trải nghiệm khó quên ha\~{}''

Iroha nghe xong thì hai hàm răng va vào nhau lập cập, mặt cắt không còn một giọt máu. Trong tâm trí non nớt của cô bé samurai, hình ảnh của tôi lúc này chẳng khác nào một ma vương tối thượng ngụy trang dưới lốt nhân viên văn phòng, một kẻ có thể búng tay đưa cô nàng lưu đày sang chiều không gian khác bất cứ lúc nào!

``S-Sống sót qua Hư Vô\dots\ phá vỡ không gian số\dots\ thanh tẩy nguyền rủa ngàn năm\dots\ K-Kẻ mạnh nhất vũ trụ là có thật sao, de gozaru!?''

Đúng ngay khoảnh khắc căng thẳng tột cùng đó, tiếng bước chân quen thuộc ở bên ngoài, tôi thản nhiên đẩy cửa bước vào văn phòng, khiến cho Iroha giật nảy mình bắn cả người lên vì ám ảnh.

``Được rồi, tranh thủ thời gian nghỉ ngơi, bây giờ chúng ta sẽ bắt đầu\dots''

Chưa để tôi dứt câu thông báo, nàng samurai đã lao phắt người xuống, cắm cúi quỳ rạp cúi gằm mặt liên tục sát đất trước mặt tôi hệt như một chiếc máy dập bản năng. ``K-Kuon-dono! Em lạy anh! Em thề từ nay về sau sẽ ngoan ngoãn, tuyệt đối không dám bén mảng làm phiền anh nữa đâu ạ! Xin đại nhân hãy tha mạng cho em qua kiếp nạn này với!''

Tôi đứng hình toàn tập, mồ hôi hột túa ra như thác lũ: ``Ơ kìa\dots\ Này, em bị làm gì mà tự dưng lại quỳ lạy dữ dội thế?''

Trên bàn, chiếc đồng hồ của tôi nhấp nháy đèn tín hiệu biểu cảm cười khẩy vô tội, trong khi các cô gái đứng xung quanh thì che miệng khúc khích cười mỉm. Ryuuhou đứng góc phòng ôm bụng cười sặc sụa: ``Đấy, đã bảo chúa công nhà em uy quyền xuyên không gian lắm cơ mà! Thức thời quỳ lạy van xin sớm là khôn ngoan đấy, Iroha-san ơi!''

``Chuyện quái gì đang xảy ra với em ấy vậy trời\dots'' tôi thở dài sườn sượt, bất lực đưa tay vò đầu bứt tai khi nghĩ đến viễn cảnh sắp tới phải tiếp nhận thêm bốn thành viên báo đời nữa của nhánh mới chuẩn bị gia nhập.

Xem ra cái cuộc đời làm quản lý của tôi từ nay về sau còn gian nan và dài dài lắm đây\dots

\end{document}